% ASP.NET Core section

% 1 – What is ASP.NET Core?
\QuestionSlide[\CategoryBadge[AspNetColor!20]{ASP.NET Core}]{* What is ASP.NET Core?}

\begin{frame}[fragile]
  \frametitle{\AnswerTitle[\CategoryBadge[AspNetColor!20]{ASP.NET Core}]{What is ASP.NET Core?}}

  {\footnotesize
  \textbf{ASP.NET Core} is Microsoft's open-source, cross-platform web framework built on .NET. It unifies MVC, Web API, and Razor Pages into one pipeline running on Kestrel, with built-in DI, a modular middleware system, and Minimal API support introduced in .NET 6.
  }
\end{frame}

% 2 – Middleware Pipeline
\QuestionSlide[\CategoryBadge[AspNetColor!20]{ASP.NET Core}]{** What is the ASP.NET Core middleware pipeline?}

\begin{frame}[fragile]
  \frametitle{\AnswerTitle[\CategoryBadge[AspNetColor!20]{ASP.NET Core}]{What is the middleware pipeline?}}

  {\footnotesize
  Middleware components form a \textbf{chain of responsibility} around the HTTP request/response. Each calls \texttt{next()} to pass control forward. Components registered first wrap components registered later, so \textbf{order matters}: exception handling should be outermost, routing and auth innermost.
  }

  \begin{minted}[fontsize=\scriptsize]{csharp}
var app = builder.Build();

app.UseExceptionHandler("/error"); // catches unhandled exceptions
app.UseHttpsRedirection();
app.UseAuthentication();           // must precede UseAuthorization
app.UseAuthorization();
app.MapControllers();              // terminal middleware

app.Run();
  \end{minted}
\end{frame}

% 3 – DI Lifetimes
\QuestionSlide[\CategoryBadge[AspNetColor!20]{ASP.NET Core}]{** What are the three DI service lifetimes in ASP.NET Core?}

\begin{frame}[fragile]
  \frametitle{\AnswerTitle[\CategoryBadge[AspNetColor!20]{ASP.NET Core}]{What are the three DI lifetimes?}}

  {\footnotesize
  \begin{itemize}
    \item \textbf{Transient} – new instance on every injection; ideal for stateless, lightweight services.
    \item \textbf{Scoped} – one instance per HTTP request; default for \texttt{DbContext}.
    \item \textbf{Singleton} – one instance for the app lifetime; use for caches or config readers.
  \end{itemize}
  Injecting Scoped into Singleton is a \textbf{captive dependency} bug; \texttt{ValidateScopes()} catches it.
  }

  \begin{minted}[fontsize=\scriptsize]{csharp}
builder.Services.AddTransient<IEmailSender, SmtpSender>();
builder.Services.AddScoped<IOrderRepository, EfOrderRepo>();
builder.Services.AddSingleton<IExchangeRateCache, ExchangeRateCache>();
  \end{minted}
\end{frame}

% 4 – Minimal APIs
\QuestionSlide[\CategoryBadge[AspNetColor!20]{ASP.NET Core}]{** What are Minimal APIs in ASP.NET Core?}

\begin{frame}[fragile]
  \frametitle{\AnswerTitle[\CategoryBadge[AspNetColor!20]{ASP.NET Core}]{What are Minimal APIs?}}

  {\footnotesize
  \textbf{Minimal APIs} (.NET 6+) define HTTP endpoints as lambda delegates on \texttt{WebApplication} — no controller classes or attributes needed. They support DI, endpoint filters, route groups, and OpenAPI while drastically reducing boilerplate.
  }

  \begin{minted}[fontsize=\scriptsize]{csharp}
var builder = WebApplication.CreateBuilder(args);
builder.Services.AddScoped<IProductService, ProductService>();
var app = builder.Build();

app.MapGet("/products/{id:int}", async (int id, IProductService svc) =>
    await svc.GetByIdAsync(id) is { } p
        ? Results.Ok(p)
        : Results.NotFound())
   .WithName("GetProduct")
   .WithOpenApi();

app.Run();
  \end{minted}
\end{frame}

% 5 – Action Filters
\QuestionSlide[\CategoryBadge[AspNetColor!20]{ASP.NET Core}]{** What are Action Filters in ASP.NET Core MVC?}

\begin{frame}[fragile]
  \frametitle{\AnswerTitle[\CategoryBadge[AspNetColor!20]{ASP.NET Core}]{What are Action Filters?}}

  {\footnotesize
  \textbf{Action Filters} are attributes that execute code before/after a controller action. They handle cross-cutting concerns — validation, logging, caching — without polluting action bodies. Pipeline order: Authorization $\to$ Resource $\to$ Action $\to$ Result $\to$ Exception.
  }

  \begin{minted}[fontsize=\scriptsize]{csharp}
public class ValidateModelFilter : ActionFilterAttribute
{
    public override void OnActionExecuting(ActionExecutingContext ctx)
    {
        if (!ctx.ModelState.IsValid)
            ctx.Result = new UnprocessableEntityObjectResult(ctx.ModelState);
    }
}

[ApiController, Route("api/[controller]")]
[ValidateModelFilter]
public class OrdersController : ControllerBase { }
  \end{minted}
\end{frame}

% 6 – Output Caching
\QuestionSlide[\CategoryBadge[AspNetColor!20]{ASP.NET Core}]{** What is Output Caching in ASP.NET Core 7+?}

\begin{frame}[fragile]
  \frametitle{\AnswerTitle[\CategoryBadge[AspNetColor!20]{ASP.NET Core}]{What is Output Caching?}}

  {\footnotesize
  \textbf{Output Caching} stores the full HTTP response and replays it without re-executing the endpoint. Configurable by duration, query-string keys, and custom policies. Backed by in-memory or distributed Redis storage. Distinct from Response Caching which relies on client cache headers.
  }

  \begin{minted}[fontsize=\scriptsize]{csharp}
builder.Services.AddOutputCache(opts =>
    opts.AddPolicy("products", p =>
        p.Expire(TimeSpan.FromMinutes(5))
         .SetVaryByQuery("category")));

app.UseOutputCache();

app.MapGet("/products", GetProducts).CacheOutput("products");
  \end{minted}
\end{frame}

% 7 – Health Checks
\QuestionSlide[\CategoryBadge[AspNetColor!20]{ASP.NET Core}]{** What are Health Checks in ASP.NET Core?}

\begin{frame}[fragile]
  \frametitle{\AnswerTitle[\CategoryBadge[AspNetColor!20]{ASP.NET Core}]{What are Health Checks?}}

  {\footnotesize
  \textbf{Health Checks} expose a \texttt{/health} endpoint reporting Healthy / Degraded / Unhealthy status. Kubernetes liveness and readiness probes, Azure App Service, and load balancers consume them. Implement \texttt{IHealthCheck} for custom dependencies (queue depth, third-party API).
  }

  \begin{minted}[fontsize=\scriptsize]{csharp}
builder.Services.AddHealthChecks()
    .AddDbContextCheck<AppDbContext>()
    .AddCheck("payments-api", () =>
        paymentGateway.IsReachable()
            ? HealthCheckResult.Healthy()
            : HealthCheckResult.Degraded("Slow response"));

app.MapHealthChecks("/health");
  \end{minted}
\end{frame}

% 8 – Options Pattern
\QuestionSlide[\CategoryBadge[AspNetColor!20]{ASP.NET Core}]{* What is the Options pattern in ASP.NET Core?}

\begin{frame}[fragile]
  \frametitle{\AnswerTitle[\CategoryBadge[AspNetColor!20]{ASP.NET Core}]{What is the Options pattern?}}

  {\footnotesize
  The \textbf{Options pattern} binds \texttt{IConfiguration} sections to strongly-typed classes via \texttt{IOptions<T>}, \texttt{IOptionsSnapshot<T>} (per-request), or \texttt{IOptionsMonitor<T>} (live reload). It provides type-safe access to configuration and supports validation with \texttt{DataAnnotations}.
  }

  \begin{minted}[fontsize=\scriptsize]{csharp}
public class EmailOptions
{
    public string Host { get; set; } = "";
    public int Port { get; set; } = 587;
}

builder.Services.Configure<EmailOptions>(
    builder.Configuration.GetSection("Email"));

// Consume
public class Mailer(IOptions<EmailOptions> opts)
{
    private readonly EmailOptions _email = opts.Value;
}
  \end{minted}
\end{frame}

% 9 – Rate Limiting
\QuestionSlide[\CategoryBadge[AspNetColor!20]{ASP.NET Core}]{*** What is Rate Limiting middleware in ASP.NET Core 7+?}

\begin{frame}[fragile]
  \frametitle{\AnswerTitle[\CategoryBadge[AspNetColor!20]{ASP.NET Core}]{What is Rate Limiting middleware?}}

  {\footnotesize
  Built-in rate limiting (ASP.NET Core 7) protects APIs via \textbf{Fixed Window}, \textbf{Sliding Window}, \textbf{Token Bucket}, or \textbf{Concurrency} limiters. Apply globally or per-endpoint by policy name. Returns HTTP 429 when limits are exceeded.
  }

  \begin{minted}[fontsize=\scriptsize]{csharp}
builder.Services.AddRateLimiter(opts =>
{
    opts.AddFixedWindowLimiter("api", limiter =>
    {
        limiter.Window = TimeSpan.FromSeconds(10);
        limiter.PermitLimit = 100;
        limiter.QueueLimit = 5;
    });
    opts.RejectionStatusCode = StatusCodes.Status429TooManyRequests;
});

app.UseRateLimiter();
app.MapGet("/api/data", Handler).RequireRateLimiting("api");
  \end{minted}
\end{frame}
